\section{软件设计与核心算法}

\subsection{软件总体结构}
\subsubsection{F407 侧}
F407 工程基于 Keil（ARMCC5）+ HAL 库，从开发板厂商的 FSMC-TFTLCD 例程裁剪而来，
调度使用 FreeRTOS-Kernel V10.6.2，直接调用内核 API，不套 cmsis\_os。

FreeRTOS 在这里解决的是周期互相拖累的问题。界面要每 5\,ms 刷新，计步要每 20\,ms
读传感器，测温和闹钟可以 200\,ms 做一次，语音上传却会在串口上停等最多 300\,ms。
这些工作若挤在同一个主循环里，一次停等就会让走时和计步一起停住。
所以按周期拆成任务：界面任务优先级最高，而且是唯一允许调用 LVGL 的任务；
链路、传感和慢任务只向一个 16 项队列投事件，不直接改控件。
链路任务优先级次之，停等发生在它自己的时间片里。
传感、慢任务、灯阵和游戏节拍优先级相同，各自用 \texttt{vTaskDelayUntil}
对齐周期。游戏开始时再创建步进任务，退出就删除，三个游戏不必同时占栈。
游戏步进只更新逻辑，画面仍回到界面任务里画。

节拍用 SysTick，1\,ms。HAL 的毫秒时基改到 TIM6，避免和调度器抢同一个滴答。
中断优先级分组取 4 位全抢占，这是该内核版本的要求。
栈深单位是字。分配失败或栈溢出时直接停机，方便在调试器里看现场。
常驻任务六个，划分如下。

\begin{table}[htbp]
\centering\small
\caption{F407 任务划分}\label{tab:tasks}
\begin{tabular}{@{}llp{2.15cm}p{7.6cm}@{}}
\toprule
任务 & 周期 & 优先级/栈 & 职责\\
\midrule
ui & 5\,ms & 3 / 1024 & 出队刷新 LVGL，唯一允许调用 lv\_* 的任务\\
link & 5\,ms & 2 / 512 & 语音状态机、C3 帧收发与心跳、秒表和倒计时\\
sensor & 20\,ms & 1 / 512 & MPU6050 采样、计步与跌倒；每 10 步请求落盘\\
slow & 200\,ms & 1 / 512 & RTC、DS18B20、闹钟、蜂鸣 1\,s 关断、EEPROM 写入\\
gtick & 随游戏 & 1 / 128 & 按 40 / 400 / 700\,ms 唤醒对应游戏步进\\
game & 进入时创建 & 1 / 256 & 贪吃蛇、乒乓或俄罗斯方块走一步，退出即删除\\
rgb & 约 90\,ms & 1 / 256 & PE5 上 5$\times$5 WS2812 点阵循环滚动 CQUT\\
\bottomrule
\end{tabular}
\end{table}

初始化按固定顺序进行：168\,MHz 时钟、调试串口、LED 与蜂鸣的安全态、
FSMC（LCD 与 SRAM）、LVGL 与界面、传感器和 EEPROM 恢复、RTC、DS18B20、
UART4，最后创建任务并启动调度器。初始化期间 LED1 兼作阶段指示。
USART2/3 以及板载其他串口不初始化，避免复用冲突。

\subsubsection{C3 侧}
C3 使用 ESP-IDF v5.4.4。app\_main 仅做 NVS 初始化与版本横幅（固件版本、编译时间、
ELF SHA-256 前 16 位、运行分区、空闲堆），随后启动三个功能模块：
link（UART1 帧收发任务，优先级 5）、asr（音频接收/上传状态机与后台预擦除任务）、
wifi\_mgr（事件驱动连接管理）。weather 为纯函数式模块。密钥与 SSID 全部隔离在
secrets.h，不入库也不进报告。

\subsubsection{版本自检机制}
两块板分处两个工具链，板上运行的镜像与磁盘源码可能不一致，这是实板调试中的真实坑。
为此 C3 启动横幅打印 ELF SHA-256 前 16 位，F407 侧留有 printf 镜像区供调试器读取；
固件发布时以 SHA-256 哈希为唯一身份。该机制在 5.4 节的排查中直接发挥了作用。

\subsection{核心算法一：音频采集与可靠传输}
\subsubsection{\texorpdfstring{16\,kHz DMA 采集}{16 kHz DMA 采集}}
录音链路为``TIM3 定时触发 + ADC1 转换 + DMA 双半环搬运''：
TIM3 由 84\,MHz 定时器时钟分频 $5250$ 得 16\,kHz TRGO，触发 ADC1 通道 5（PA5）
以 84 周期采样时间转换；DMA 使用 2048 点环形缓冲，分为两个 1024 点半块，
半满/全满中断中把当前半块转存为 16\,bit PCM：
\[
\text{sample} = (adc - 2048) \times 16
\]
即把 12\,bit ADC 的中点搬到 0 并左移 4 位扩展为 16\,bit。
DMA 在填一个半块的同时 ISR 拷贝另一个半块，实现无丢失连续采集。
松开按键时按冻结的 NDTR 寄存器值补齐尾部不满半块。

录音缓冲不放在内部 SRAM：按 10\,s 上限需要 320\,000\,B（16\,kHz$\times$2\,B$\times$10\,s），
故 PCM 缓冲经分散加载放入外部 SRAM（0x68000000，占其 1\,MB 的 30.5\%），
内部 RAM 只保留 DMA 环形区与句柄。考虑到外部 SRAM 挂在 FSMC 上存在时序风险，
每个半字写入后立即读回核对，失败置 MIC\_ERR\_MEMORY 错误并禁止上传——
这个自检在联调中拦截过真实的 SRAM 时序问题（见 5.4 节）。

\subsubsection{分包停等传输}
上传把 PCM 划分为不超过 506\,B 的段。偏移 4\,B 加数据 506\,B 构成 510\,B 载荷；
帧头 9\,B（SOF 到 LEN）加 CRC 2\,B，音频数据帧共 521\,B。
满长 10\,s 音频为 BEGIN + 633 个 DATA + END，共 635 帧。
传输采用停等式流控：每帧发出后等待对端回显 SEQ 的 ACK，超时 300\,ms 重发原帧，
最多 3 次；对端 NACK 或 3 次超时则本次会话失败。链路层统计 crc/ore/drop/stale\_ack
四类计数并周期打印，作为通信质量指标（见 5.2 节）。

\subsection{核心算法二：计步与跌倒粗判}
计步不使用 MPU6050 的 DMP 硬件运动处理单元，算法完全自研、参数可标定可解释。
MPU6050 配置为 $\pm$2\,g 量程、50\,Hz 输出（SMPLRT\_DIV=19、DLPF$\approx$44\,Hz），
每 20\,ms 读三轴原始值，流程：
\begin{align}
mag &= \sqrt{(a_x/16384)^2+(a_y/16384)^2+(a_z/16384)^2}\\
base &\leftarrow 0.98\,base + 0.02\,mag \quad (\text{初值 }1\,g)\\
dyn &= mag - base
\end{align}
一阶滤波以慢时间常数跟踪重力与身体姿态构成的基线，动态分量 dyn 反映脚步冲击。
记步条件：$dyn$ 上穿阈值 $th$（默认 150\,mg）后回落，且回落时刻距上一步不小于
gap（默认 300\,ms）。两个参数均可通过校准页现场调整（阈值 50--500\,mg、间隔四档），
出厂值针对正常步幅与行走频率设定。

跌倒粗判叠加在同一数据流上：合加速度 $>2.0\,g$ 连续两个采样（冲击），
或 $<0.4\,g$ 持续 80\,ms（自由落体失重）后 1\,s 内出现冲击，即判定跌倒；
触发后 5\,s 静默抑制重复告警。判定成立时蜂鸣 1\,s、切换告警页。
该算法为演示性粗判，不是医疗器械级别的跌倒检测。

\subsection{核心算法三：自定义帧协议}
F407 与 C3 之间定义了一种轻量可靠帧，双方各自独立实现（C3 侧为 11 态字节流状态机），
帧结构见表~\ref{tab:frame}。

\begin{table}[htbp]
\centering\small
\caption{帧结构与字段}\label{tab:frame}
\begin{tabular}{@{}clp{7.2cm}@{}}
\toprule
字段 & 长度 & 说明\\
\midrule
SOF & 2 & 0xA5 0x5A，字节流重新同步标志\\
VER & 1 & 0x01\\
SRC/DST & 各 1 & 0xF4（F407）/ 0xC3（ESP32-C3）\\
TYPE & 1 & 帧类型，见表~\ref{tab:types}\\
SEQ & 1 & 发送方递增；应答帧回显被应答帧的 SEQ\\
LEN & 2 & 载荷字节数（小端），0--512\\
PAYLOAD & LEN & 多字节整数一律小端\\
CRC & 2 & Modbus CRC16（多项式 0xA001、初值 0xFFFF），覆盖 VER 到载荷末字节（不含 SOF）\\
\bottomrule
\end{tabular}
\end{table}

\begin{table}[htbp]
\centering\small
\caption{帧类型定义}\label{tab:types}
\begin{tabular}{@{}llp{6.2cm}@{}}
\toprule
TYPE & 名称与方向 & 载荷\\
\midrule
0x01 & HELLO（F4$\to$C3） & 空；上电后每 2\,s 一次直到获 ACK\\
0x02 & HELLO\_ACK（C3$\to$F4） & 1\,B 在线标志（后接 4\,B IPv4）\\
0x10 & AUDIO\_BEGIN（F4$\to$C3） & PCM 总字节数 u32 小端\\
0x11 & AUDIO\_DATA（F4$\to$C3） & u32 偏移 + $\le$506\,B PCM\\
0x12 & AUDIO\_END（F4$\to$C3） & 空；C3 开始上传云端\\
0x13 & ASR\_TEXT（C3$\to$F4） & UTF-8 文本，长度由 LEN 给出\\
0x20 & CMD\_WEATHER（F4$\to$C3） & 城市名 UTF-8 字节，空则用默认城市\\
0x21 & WEATHER\_RESULT（C3$\to$F4） & 一行中文摘要 UTF-8\\
0x30 & ALERT\_FALL（F4$\to$C3） & 告警事件上报\\
0x7E/0x7F & ACK / NACK & NACK 载荷 1\,B 错误码（0x01 超时 / 0x02 CRC / 0x03 未联网 / 0x04 HTTP / 0x05 空结果 / 0x06 忙）\\
\bottomrule
\end{tabular}
\end{table}

接收侧逐字节喂入状态机：搜 SOF、核对长度合法性、收满载荷后校验 CRC，
错帧静默丢弃并计数（不回 NACK，避免把错误帧放大成双向风暴）；
语音音频帧启用停等重传，其余帧为尽力而为。
设计要点是\textbf{把可靠性做在协议里}：460800\,bit/s、8N1 下，521\,B 帧的线上时间约 11\,ms。
停等窗口因此很小；重传 3 次仍失败，就应当停下来查线，而不是继续空转。

\subsection{云端接入（C3 侧）}
\subsubsection{语音上传}
C3 把 F407 送来的三段音频帧重组、暂存到专用 flash audio 分区（384\,KB，
空闲时后台预擦除扇区，接收时 4\,KB 缓冲写入）。F407 发送 AUDIO\_END 后，
C3 从分区分块读出 PCM，边生成 44\,B 标准 WAV 头、边做 Base64 流式编码并写入
HTTPS 请求体——任何时刻都不在 RAM 中保留整段 PCM/WAV/Base64/JSON，
10\,s 音频（320\,KB）在 C3 上传输的峰值内存保持在 KB 级。
请求发往火山方舟语音模型接口（模型 doubao-seed-2-0-lite，密钥不入库），
提示词约束模型``直接回应用户语音内容、回答不超过 40 个汉字''以适配手表显示。
云端返回的 JSON 经解析提取文本字段，以 ASR\_TEXT 帧回传 F407 原样显示，
F407 不再做关键词解析。失败链路区分连接失败、上传失败、HTTP 非 2xx、JSON 解析
失败四级日志，便于定位问题所在层。

\subsubsection{天气查询}
天气服务选用免密钥的 Open-Meteo。CMD\_WEATHER 帧允许携带中文城市名：
C3 先调用 Open-Meteo 地理编码接口把城市名解析为经纬度，
再请求 forecast 接口（当前温度、体感温度、湿度、风向、天气代码），
把 WMO 天气代码映射为中文（晴/多云/阴/雾/雨/雪/雷暴等），
拼装成一行固定顺序的摘要``\,城市\ 现象\ 温度$^\circ$C\ 体感\ 湿度\%\ 风向\,''回传。
默认城市写在固件配置里，F407 可循环切换重庆、北京、上海、广州、成都。

\subsubsection{Wi-Fi 管理}
wifi\_mgr 为事件驱动：掉线后每 5\,s 重连，状态变化即时以 HELLO\_ACK(在线标志)
通知 F407，由 F407 更新状态栏的``在线/离线''显示。语音会话期间暂缓天气请求，
避免两个 HTTPS 会话并发竞争内存。

\subsection{GUI 实现要点}
LVGL 8.3 使用 20 行双缓冲，每块 320$\times$20，只刷新脏矩形。
flush 经 16 位 FSMC 写入 HX8357D 的 GRAM，不建全屏帧缓冲。
界面任务每 5\,ms 调用 lv\_timer\_handler。
主页六行卡片的焦点用 150\,ms 缓动切换。按键语义按页面分发，新增页面不必改调度循环。
当前镜像里的汉字字体是 2.6 节所说的 trial 字库，不是脚本生成的那份 KB 级子集。

\subsection{系统部署}
\subsubsection{构建与烧录流程}
\begin{itemize}
  \item F407：EIDE / Keil ARMCC5，产物 Template.hex。当前镜像代码约 245\,KB，
        Flash 执行区约 836\,KB / 1\,MB，片内 RAM 约 105.5\,KB / 128\,KB，
        外扩 SRAM 中的录音缓冲 320\,000\,B。OpenOCD 0.12 + CMSIS-DAP，
        SWD 500\,kHz 烧录并核对。
  \item C3：ESP-IDF PowerShell 终端 idf.py 构建，产物 ai-watch-c3.bin；
        esptool 先以 chip-id 只读核对芯片身份再烧录，串口监视 115\,200\,bps。
  \item 版本核对：烧录后以启动横幅的 SHA-256 与本地哈希比对确认板上即磁盘版本。
\end{itemize}

\subsubsection{运行环境}
系统上电自启，无需调试器介入：F407 从 Flash 启动（BOOT0 接地），
C3 从其 flash 启动并自动连 Wi-Fi；两板通过 UART4 双绞/杜邦线连接。
演示仅需 5\,V 电源与一块已写入 Wi-Fi 凭据的 C3。
